\documentclass{article}

\usepackage{amsmath,amssymb}
\usepackage{xurl}
\usepackage[colorlinks=true,linkcolor=blue,citecolor=blue,urlcolor=blue]{hyperref}

% Shared commands for notes in docs/paper-gaps/.

\DeclareUnicodeCharacter{2080}{\ifmmode{}_{0}\else\textsubscript{0}\fi}
\DeclareUnicodeCharacter{2081}{\ifmmode{}_{1}\else\textsubscript{1}\fi}
\DeclareUnicodeCharacter{2082}{\ifmmode{}_{2}\else\textsubscript{2}\fi}
\DeclareUnicodeCharacter{2099}{\ifmmode{}_{n}\else\textsubscript{n}\fi}

\newcommand{\Lean}{\textsc{Lean}}
\ExplSyntaxOn
\NewDocumentCommand{\leanid}{m}{
  \ifmmode
    \textsf{\detokenize{#1}}
  \else
    \group_begin:
    \urlstyle{sf}
    \tl_set:Nn \l_tmpa_tl {#1}
    \tl_replace_all:Nnn \l_tmpa_tl {\_} {_}
    \exp_args:NV \path \l_tmpa_tl
    \group_end:
  \fi
}
\ExplSyntaxOff
\providecommand{\tr}{\operatorname{tr}}
\newcommand{\tnleanIssueBase}{https://github.com/LionSR/TNLean/issues/}
\newcommand{\tnleanPrBase}{https://github.com/LionSR/TNLean/pull/}
\newcommand{\tnleanDocBase}{https://sirui-lu.com/TNLean/docs/}
\newcommand{\ghissue}[1]{\href{\tnleanIssueBase#1}{GitHub issue \##1}}
\newcommand{\ghpr}[1]{\href{\tnleanPrBase#1}{PR \##1}}
\newcommand{\doclink}[1]{\href{\tnleanDocBase#1.html}{\path{#1}}}

% Dirac notation and the identity operator, as in blueprint/src/macros/common.tex.
% \providecommand keeps note-local definitions authoritative.
\usepackage{dsfont}
\providecommand{\ket}[1]{|#1\rangle}
\providecommand{\bra}[1]{\langle #1|}
\providecommand{\braket}[2]{\langle #1 | #2 \rangle}
\providecommand{\ketbra}[2]{|#1\rangle\!\langle #2|}
\providecommand{\Id}{\mathds{1}}
\providecommand{\one}{\mathds{1}}
\providecommand{\C}{\mathbb{C}}
\providecommand{\MN}[1]{M_{#1}(\C)}
\providecommand{\MD}{M_D(\C)}

% External referencing for paper sources.  The build helper writes sanitized
% auxiliary files under build/paper-aux/ and excludes duplicated source labels.
\usepackage{xr}

% Import a generated paper auxiliary file with a caller-supplied label prefix.
% The candidate paths support builds from docs/paper-gaps/, the repository root,
% and build/paper-aux/ itself.
\newcommand{\importpaperaux}[2]{%
  \IfFileExists{../../build/paper-aux/#2.aux}{%
    \externaldocument[#1]{../../build/paper-aux/#2}%
  }{%
    \IfFileExists{build/paper-aux/#2.aux}{%
      \externaldocument[#1]{build/paper-aux/#2}%
    }{%
      \IfFileExists{#2.aux}{%
        \externaldocument[#1]{#2}%
      }{}%
    }%
  }%
}

% General external reference macros
\newcommand{\paperref}[2]{\ref{#1-#2}}
\newcommand{\papereqref}[2]{(\ref{#1-#2})}

% CPSV16 source (1606.00608 / MPDO-22-12-17-2.tex) external document and shortcuts
\importpaperaux{cpsv16-}{1606.00608/MPDO-22-12-17-2-tnlean}
\newcommand{\cpsvref}[1]{\paperref{cpsv16}{#1}}
\newcommand{\cpsveqref}[1]{\papereqref{cpsv16}{#1}}

% Verdict marker: \gapnote{<kind>}{<status>}. Kind is the policy
% classification (clarification, local-correction, scope-restriction,
% unfaithful, false-source, open-gap); status is open, wip (elimination
% underway), resolved, or historical. Machine-read by the site index;
% prints a verdict line.
\newcommand{\gapnote}[2]{\par\noindent\textbf{Verdict:} #1 (#2).\par}


\title{Schuch--P\'erez-Garc{\'\i}a--Cirac SPT Interpolation: Upper-Range Correction}
\author{}
\date{2026-07-26}

\begin{document}
\maketitle
\gapnote{local-correction}{resolved}

\section{Source construction}

Section~II.F of arXiv:1010.3732 connects two symmetric injective matrix product
states after replacing them by their isometric forms. Let their bond dimensions
be $D_0$ and $D_1$. The construction embeds the two bond spaces as consecutive
direct summands of an ambient space whose dimension is
\begin{align}
    D_{\mathrm{amb}}
        &=
        D_0+D_1.
    \label{eq:spc11_spt_interpolation_upper_range_ambient_dimension}
\end{align}
It then interpolates between maximally entangled vectors supported on the two
summands.

The source display in
\path{Papers/1010.3732/paper_v3.tex}, lines~900--903, reads
\begin{align}
    \ket{\omega_{\mathrm{src}}(\gamma)}
        &=
        (1-\gamma)\sum_{i=1}^{D_0}\ket{i,i}
        +
        \gamma\sum_{i=D_0+1}^{D_1}\ket{i,i}.
    \label{eq:spc11_spt_interpolation_upper_range_source_interpolation}
\end{align}
The source immediately states that this vector has bond dimension
$D_0+D_1$.

\section{Correction}

The second summand must occupy the $D_1$ coordinates immediately following
the first $D_0$ coordinates. Its index set is
\begin{align}
    I_1
        &=
        \{D_0+1,\ldots,D_0+D_1\},
    \label{eq:spc11_spt_interpolation_upper_range_second_index_set}
\end{align}
rather than $\{D_0+1,\ldots,D_1\}$. The corrected interpolation is
\begin{align}
    \ket{\omega(\gamma)}
        &=
        (1-\gamma)\sum_{i=1}^{D_0}\ket{i,i}
        +
        \gamma\sum_{i=D_0+1}^{D_0+D_1}\ket{i,i}.
    \label{eq:spc11_spt_interpolation_upper_range_corrected_interpolation}
\end{align}

At the first endpoint,
\begin{align}
    \ket{\omega(0)}
        &=
        \sum_{i=1}^{D_0}\ket{i,i},
    \label{eq:spc11_spt_interpolation_upper_range_first_endpoint}
\end{align}
which is the maximally entangled vector on the first bond summand. At the
second endpoint,
\begin{align}
    \ket{\omega(1)}
        &=
        \sum_{i=D_0+1}^{D_0+D_1}\ket{i,i},
    \label{eq:spc11_spt_interpolation_upper_range_second_endpoint}
\end{align}
which is the maximally entangled vector on the second bond summand. The upper
limit $D_0+D_1$ agrees with the ambient dimension in
(\ref{eq:spc11_spt_interpolation_upper_range_ambient_dimension}). By contrast,
the source range can be empty when $D_0\geq D_1$.

\section{Blueprint and formal status}

The corrected formula
(\ref{eq:spc11_spt_interpolation_upper_range_corrected_interpolation}) occurs in
the proof of the blueprint theorem \path{thm:spt_classification} in
\path{blueprint/src/chapter/ch12_symmetry_spt.tex}. An inline
\path{Local fix} comment there points to this note.

The separation argument at source lines~929--953 is explicitly restricted to
smooth MPS paths. The blueprint theorem therefore concerns symmetric MPS paths
and their parent Hamiltonians. It does not claim invariance under arbitrary
gapped Hamiltonian paths.

The formal predicate \leanid{MPSTensor.IsSameSPTPhase} records equality of the
virtual cocycle labels; it does not construct the interpolation. The blueprint
classification theorem remains marked \path{notready}. No \Lean{} declaration
implements this path or depends on the corrected summation range.

\section{Verdict}

Replacing the upper limit $D_1$ in
(\ref{eq:spc11_spt_interpolation_upper_range_source_interpolation}) by
$D_0+D_1$, as in
(\ref{eq:spc11_spt_interpolation_upper_range_corrected_interpolation}), is a
local source correction. It restores the intended direct-sum coordinates and
is the version used in the blueprint. Restricting the classification statement
to MPS and parent-Hamiltonian paths is a scope clarification that matches the
source proof.

\end{document}
